\section{Method}
\label{sec:methods}

\subsection{A mixture of catalog redshift priors}
\label{sec:methods:mixture}

We use the standard dark-siren population likelihood
\citep{ChenFishbachHolz2018, MandelFarrGair2019, Gray2020}. For event $i$ with
data $d_i$ and parameters $\Lambda$ containing the cosmology and the source
population, the single-event evidence is
\begin{equation}
  Z_i(\Lambda) = \int \dd\theta\; p(d_i \,|\, \theta)\,
                 p_{\rm pop}(m_1, m_2, \chi \,|\, \Lambda)\,
                 p(z, \hat{n} \,|\, \Lambda),
  \label{eq:zi}
\end{equation}
where $\theta$ collects source-frame masses, effective spin, redshift and sky
direction, and $p(z, \hat{n}|\Lambda)$ is the line-of-sight redshift prior
supplied by the catalog.

When more than one catalog traces the same structure, that prior becomes a
mixture,
\begin{equation}
  p(z, \hat{n} \,|\, \Lambda) = \sum_{k=1}^{K} f_k \,
      p_k(z, \hat{n} \,|\, \Lambda), \qquad \sum_k f_k = 1,
  \label{eq:mixture}
\end{equation}
with $f_k$ the fraction of mergers hosted by the population that catalog $k$
traces. We take $K = 2$, with $k = 1$ the galaxies and $k = 2$ the AGN, so
$f_2 \equiv \fagn$ and $f_1 = 1 - \fagn$. Written this way \fagn\ is a mixture
weight and nothing else: it does not enter the mass or spin distributions, and
it is identified entirely by which catalog's redshift and sky structure the
events prefer.

\subsection{The redshift prior a catalog supplies}
\label{sec:methods:prior}

Each catalog's prior is a kernel-density estimate in redshift over the
catalogued objects in the event's HEALPix pixel \citep{Gorski2005, Gray2022},
\begin{equation}
  p_k(z, \hat{n} \,|\, \Lambda) \;\propto\; g(z; \Lambda)
     \sum_{j \,\in\, k,\; \hat{n}_j = \hat{n}}
     \frac{w_j}{W_k}\,
     \frac{\mathcal{N}(z; z_j, \sigma_j)}{\zeta_j},
  \label{eq:pcat}
\end{equation}
where $z_j$, $\sigma_j$ and $w_j$ are a catalogued object's redshift, redshift
uncertainty and weight, $g(z;\Lambda)$ is the redshift weighting set by the
comoving volume element and the merger-rate evolution, and
$\zeta_j = \int \mathcal{N}(z; z_j, \sigma_j)\, g(z;\Lambda)\, \dd z$
normalises each kernel against that weighting.

The normalisation $W_k$ is what makes \fagn\ identifiable. We normalise each
catalog's sky weighting over the survey as a whole, $W_k = \sum_j w_j$ summed
over the entire catalog, rather than over the objects in each pixel separately.
The number of catalogued hosts in a pixel then survives into the likelihood, so
a pixel rich in galaxies but poor in AGN contributes weight to the first
component of Equation~(\ref{eq:mixture}) and little to the second. Under a
per-pixel normalisation every occupied pixel is rescaled to unit weight and the
sky-density contrast between the two catalogs (the difference in where and
in how strongly each traces the density field) is discarded.

The kernel widths $\sigma_j$ must stay well below the redshift resolution of
the distance measurements,
\begin{equation}
  \sigma_z^{\rm PE} \;\simeq\; \sigma_{d_L}\, d_L \left(
     \frac{\dd d_L}{\dd z} \right)^{-1},
  \label{eq:sigmazpe}
\end{equation}
or the redshift information in the likelihood comes from the overlap of a broad
prior with the distance posterior rather than from the catalogued hosts. For a
real survey this is a requirement on photometric redshifts. In the catalogs of
\S\ref{sec:data} the kernel widths are a factor \KdePeRatio\ below
$\sigma_z^{\rm PE}$ at the median detected redshift.

\subsection{Incomplete catalogs}
\label{sec:methods:completion}

A flux-limited survey stops finding hosts beyond some redshift, and the
analysis must say what is missing. We model the missing objects of each catalog
with a smooth comoving density,
\begin{equation}
  \frac{\dd N_{{\rm miss},k}}{\dd z} =
    n_{0,k}\, \frac{\dd V_c}{\dd z}\, (1+z)^{\delta_k}
    - \frac{\dd N_{{\rm obs},k}}{\dd z},
  \label{eq:miss}
\end{equation}
distributed isotropically over the sky, so the completeness is not a free
function but follows from the assumed density and the observed counts,
\begin{equation}
  C_k(z) = \frac{\dd N_{{\rm obs},k}/\dd z}
                {n_{0,k}\, (\dd V_c/\dd z)\, (1+z)^{\delta_k}}.
  \label{eq:completeness}
\end{equation}
Each catalog's prior is then Equation~(\ref{eq:pcat}) plus a term carrying
$\dd N_{{\rm miss},k}/\dd z$ spread uniformly over the sky, and the mixture of
Equation~(\ref{eq:mixture}) runs over those completed priors. Setting
$n_{0,k} \rightarrow 0$ removes the missing budget and recovers the
complete-catalog case.

This is the mechanism by which a shallow survey still constrains \fagn: as
catalogued hosts disappear each prior migrates from its own kernel-density
estimate onto its own smooth density, and the two smooth densities still differ
by the number-density ratio. The pair $(n_{0,k}, \delta_k)$ therefore fixes how
much of each population the analysis believes it cannot see, and \nzagn\ is the
one that matters: because \fagn\ is a weight between the two priors, an error
in \nzagn\ moves the AGN prior's out-of-catalog mass and \fagn\ follows. In
the analyses of this paper the catalogs are complete and the missing budget is
switched off; the flux-limited analyses of \S\ref{sec:results:incomplete}
anchor $(n_{0,k}, \delta_k)$ at the simulation's own comoving densities.

\subsection{The selection function}
\label{sec:methods:selection}

Over $\Nobs$ detected events the likelihood is
\begin{equation}
  \ln \mathcal{L}(\Lambda) = \sum_{i=1}^{\Nobs} \ln Z_i(\Lambda)
      - \Nobs \ln \mu(\Lambda),
  \label{eq:like}
\end{equation}
with
\begin{equation}
  \mu(\Lambda) = \int \dd\theta\; P_{\rm det}(\theta)\, p(\theta \,|\, \Lambda)
  \label{eq:mu}
\end{equation}
the detectable fraction \citep{MandelFarrGair2019, Vitale2022}. Because
Equation~(\ref{eq:zi}) is a density over the observed data, $P_{\rm det}$ must
be a property of the observed data as well; \S\ref{sec:data:events} constructs
the simulated detection rule that way.

We estimate $\mu$ by importance sampling a set of simulated signal injections
stored in the cosmology-independent observables
$(m_1^{\rm det}, q, d_L)$ together with the density each was drawn from, and
require the effective sample size of that Monte Carlo estimate to satisfy
$\Neff > \NeffFactor\,\Nobs$ \citep{Farr2019}. Values of $\Lambda$ that fail
the criterion are not quoted.

Because $\mu$ is conditioned on the catalogs through Equation~(\ref{eq:pcat}),
an injection contributes only if its redshift falls within a few $\sigma_j$ of
a catalogued object in its own pixel. Injections drawn from the source
population almost never land on the sparse catalog, so at large \fagn\ the
estimate of $\mu$ is dominated by a handful of samples and the recovered
distance scale tilts. We therefore draw injections from a three-component
mixture: from the source population, from a distribution uniform in comoving
volume, and from a component placed on the catalogued AGN themselves. Each
injection stores the exact mixture density at its own coordinates, so the
estimate of Equation~(\ref{eq:mu}) remains valid for any component weights. The
third component reads the pixelated catalog file the likelihood conditions on,
so its proposal and the target cannot disagree through a pixelisation or
kernel-width mismatch. The injection sets are drawn once, against the
complete catalogs, and reused at every survey depth; the
effective-sample-size criterion is enforced at every depth.

\subsection{Evaluating the likelihood}
\label{sec:methods:inference}

At most two parameters are free at a time, so we evaluate
Equation~(\ref{eq:like}) on a grid rather than sampling it. Priors are flat, and
we quote medians with equal-tailed 68\% intervals; where the maximum lies on a
boundary of the scanned range we quote no interval. $\Omega_{m,0}$ is fixed at
\OmZeroTruth\ and the mass and spin distributions are held at their input
values, so the mass distribution is informative but not free. The free
parameters are \Hz, scanned over
$[\HzeroScanLo, \HzeroScanHi]$~\kmsmpc, and \fagn.
